\chapter{איך נבנה המסמך: צוות סוכני \en{AI}}\label{ch:how-built}

מסמך זה לא נכתב בעורך טקסט. הוא \emph{נבנה} כמו פרויקט תוכנה: ראש צוות אנושי הגדיר חזון ודרישות, וצוות של סוכני \en{AI} כתב, איֵר, בדק ותיקן. הנספח מתעד את התהליך, כי בעולם שבו כל אחד יכול להפיק \en{PDF} בלחיצה, הערך המקצועי נמצא באפיון, בהובלה ובבקרה.

\section{מבנה הצוות}

איור~\ref{fig:agents} מציג את הצוות. ה-\en{Orchestrator} (מודל \en{Opus}) בנה את התבנית, את מודל הנתונים ואת האיורים. שלושה כותבים (מודל \en{Sonnet}) כתבו את הפרקים במקביל, כל אחד בקבצים שלו בלבד. זוג סוכנים של "צוות אדום", \en{Detector} ו-\en{Fixer}, סגר את הלולאה. חלוקת המודלים היא החלטה של ניהול משאבים: המודל החזק מקבל את ההחלטות הארכיטקטוניות, שהן מעטות וקריטיות, והמודל המהיר מקבל את רוב הטקסט.

\begin{figure}[htbp]
  \centering
  \resizebox{0.86\linewidth}{!}{% How this document was produced: the agent team (appendix).
\begin{tikzpicture}[node distance=7mm and 6mm,
  ag/.style={rounded corners=2mm,draw=rrForest,line width=0.8pt,fill=#1,
             minimum width=3.1cm,minimum height=11mm,align=center,font=\scriptsize},
  ag/.default=rrCream,
  ar/.style={-{Stealth[length=2.5mm]},line width=0.9pt,color=rrForest}]
  \node[ag=rrTerra!30] (human) {\textbf{ראש צוות (אדם)}\\ חזון · אישור \en{PRD}};
  \node[ag=rrForest!20,below=of human] (orch) {\textbf{\en{Orchestrator}}\\ \en{Opus} · \en{CLS} · נתונים};
  \node[ag,below left=9mm and -4mm of orch] (w2) {\textbf{\en{writer-market}}\\ \en{Sonnet} · פרקים 4–6};
  \node[ag,left=of w2] (w3) {\textbf{\en{writer-finance}}\\ \en{Sonnet} · פרקים 7–10};
  \node[ag,right=of w2] (w1) {\textbf{\en{writer-strategy}}\\ \en{Sonnet} · פרקים 1–3};
  \node[ag=rrLeaf!40,below=14mm of w2] (det) {\textbf{\en{qa-detector}}\\ צוות אדום · 8 \en{Skills}};
  \node[ag=rrSky!60,right=12mm of det] (fix) {\textbf{\en{qa-fixer}}\\ \en{CLS} / \en{Skill} בלבד};
  \draw[ar] (human) -- (orch);
  \foreach \w in {w1,w2,w3} \draw[ar] (orch) -- (\w);
  \foreach \w in {w1,w2,w3} \draw[ar] (\w) -- (det);
  \draw[ar,rrTerra] (det) -- node[above,font=\scriptsize]{דוח} (fix);
  \draw[ar,dashed] (fix.east) -- ++(0.7,0) |- node[pos=0.25,right,font=\scriptsize]{סבב חוזר} (orch.east);
\end{tikzpicture}
}
  \caption{צוות הסוכנים שהפיק את המסמך וזרימת העבודה ביניהם}
  \label{fig:agents}
\end{figure}

\section{שלושה עקרונות הנדסיים}

\begin{enumerate}
  \item \textbf{מקור אמת יחיד.} כל מספר במסמך נגזר מקובצי \en{CSV} בתיקייה \en{data}. סקריפט \en{Python} מפיק מהם את הגרפים, את הטבלאות ואת קובץ פקודות המאקרו של המספרים. כותב שמקליד מספר ידנית במקום להשתמש במאקרו נתפס בסריקה האוטומטית.
  \item \textbf{חוק הברזל של הרפרודוקטיביליות.} תקלה מערכתית מתוקנת בתבנית (\en{roofroot.cls}, עם העלאת גרסה) או ב-\en{Skill}, ולעולם לא בשורה בודדת של טקסט. כך התיקון "מסתיר" את הלקח ומונע את הישנות התקלה.
  \item \textbf{בנייה בפקודה אחת.} \en{make} מריץ את \en{Python}, \en{LuaLaTeX} ו-\en{biber} ומפיק \en{PDF} זהה בכל מחשב. הגופנים (\en{Culmus}, \en{TeX Gyre}, \en{Fira}) מגיעים מהפצת ה-\en{TeX}, ולא מהמערכת.
\end{enumerate}

\section{משפחת ה-\en{Skills} לבקרת איכות}

כל \en{Skill} הוא קובץ הוראות שסוכן טוען לפני שהוא עובד, ומחזיק "עדשה" אחת של בדיקה. טבלה~\ref{tab:skills} מסכמת אותם.

\begin{table}[htbp]
  \centering\small
  \rrstripes
  \begin{tabularx}{\linewidth}{>{\raggedright\arraybackslash}p{3cm}Y}
    \rrtablehead \rrth{\en{Skill}} & \rrth{תפקיד} \\
    \en{qa-super}       & מתזמר את כל הבדיקות, מריץ סריקה אוטומטית ובדיקה חזותית של כל עמוד, וכותב דוח סבב \\
    \en{qa-BiDi}        & תקינות כיווניות: מונחים באנגלית, מספרים שליליים, סוגריים ותבליטים \\
    \en{qa-table}       & יישור טבלאות, סדר עמודות מימין לשמאל, ושימוש במאקרו לכל מספר \\
    \en{qa-code}        & קטעי קוד: אורך שורה, כיוון \en{LTR} והתאמה לקוד המודל \\
    \en{qa-img}         & איורים: מקור (קוד בלבד), כיתוב, הפניה לפני ההופעה וגלישת טקסט \\
    \en{qa-typeset}     & אזהרות קומפילציה: \en{Overfull}, הפניות חסרות וגליפים חסרים \\
    \en{qa-infra}       & מבנה הפרויקט ובדיקת רפרודוקטיביליות \\
    \en{qa-cls-version} & עקביות גרסת התבנית ויומן השינויים שלה \\
    \en{writer-he-en}   & מדריך סגנון לכותבים: קול, חוזה \en{LaTeX} ומקורות \\
    \en{fig-python}, \en{tikz-artist} & מוסכמות לגרפים ולאיורים בקוד \\
  \end{tabularx}
  \caption{ה-\en{Skills} של הפרויקט (הקבצים המלאים בתיקייה \en{.claude/skills} במאגר)}
  \label{tab:skills}
\end{table}

\section{חקר אירועים}

כל תקלה שהתגלתה תועדה כ"חקר אירוע" בפורמט ללא האשמה: מה קרה, למה זה קרה, למה זה לא התגלה, ומה שונה כדי שזה לא יחזור. שלוש דוגמאות מתוך שלב בניית התבנית:

\begin{riskbox}[\en{INC-001}: קוד \en{Python} הודפס מימין לשמאל]
\en{listings} בונה כל שורת קוד בקופסה נפרדת, שירשה את כיוון המסמך. התוצאה: מילים בסדר הפוך וסוגריים מהופכים. התיקון נעשה בתבנית, באמצעות וו (\en{hook}) שמכריח \en{LTR} ושפה אנגלית בתוך כל רשימת קוד, ונוספה בדיקה ל-\en{qa-BiDi}.
\end{riskbox}

\begin{riskbox}[\en{INC-002}: טבלה אוטומטית לא התקמפלה]
\en{\textbackslash rowcolor} חייב להיות האסימון הראשון בשורה. כשהשורות נטענו עם \en{\textbackslash input} בתוך \en{tabular}, הוא כבר לא היה ראשון. התיקון: המחולל ב-\en{Python} מפיק את סביבת הטבלה המלאה, ונוסף כלל ל-\en{qa-table}.
\end{riskbox}

\begin{riskbox}[\en{INC-003}: סימן המינוס קפץ לצד הלא נכון]
בטבלת רווח והפסד, \en{-1,413} הוצג כ-"1,413-" בהקשר עברי. המחולל עוטף מעתה כל מספר ב-\en{LTR}, ונוספה בדיקה אוטומטית למספרים שליליים חשופים בטקסט.
\end{riskbox}

\begin{insightbox}[הלקח הניהולי]
הסוכנים כותבים מהר. היתרון האנושי הוא לראות את מה שהם לא רואים, ולהפוך כל תקלה לכלל שנאכף אוטומטית. דוחות סבבי הצוות האדום, קובצי הפרומפטים וחקרי האירועים המלאים נמצאים במאגר ה-\en{GitHub} של הפרויקט.
\end{insightbox}
